\origpage{339}{N. Robertson, P.D. Seymour / Journal of Combinatorial Theory, Series B 90 (2004) 325--385}
If for some $i \ge 1$ there is a component $K'$ of $K$ with
\begin{equation*}
|V(K') \cap V(\Gamma_i)| + |V(K') \cap bd(\Sigma_\chi)| \le 1,
\end{equation*}
then $V(K') \cap bd(\Sigma_\chi) = \emptyset$, and so $K' = K = F$, and $|F \cap V(\Gamma_i)| \le 1$, which is impossible since $\Gamma_i$ is internally 3-connected. Thus 5.3 can be applied. Now by 5.4, $U(K)$ meets at most one cuff (for $K$ is a subgraph of each $\Gamma_i^*$, and yet $dist(\Gamma_i) > \frac{1}{2}|E(F)|$ for all sufficiently large $i$). Moreover, $|U(K) \cap bd(\Sigma_\chi)| \le 2$, with equality only if $K \ne F$ (from our choice of the $F_i$'s.) In particular, when we cut along $U(K)$ (with surjection $\psi$) we obtain either one or two fragments. For each fragment $\Sigma_1$ say, $\hat{\Sigma}_1$ is simpler than $\hat{\Sigma}_\chi$ because $U(F)$ is non-null-homotopic in $\hat{\Sigma}$, and so if there is only one fragment the proof is completed as for 5.4, using \textbf{S}$_1$ in place of \textbf{S}$_3$. (Since we are applying \textbf{S}$_1$ it does not matter whether the fragment is orientable.) Suppose then that there are two fragments $\Sigma_1, \Sigma_2$. For $t = 1, 2$, let $\chi_t$ be the colour scheme defined by $\Sigma_{\chi_t} = \Sigma_t$, $\Phi_{\chi_t} = \psi^{-1}(\Phi) \cap \Sigma_t$ etc., as in 5.4. For each $i \ge 1$, let $\Gamma_i' = \psi^{-1}(\Gamma_i) \cap \Sigma_1$, $\Gamma_i'' = \psi^{-1}(\Gamma_i) \cap \Sigma_2$. For $e \in E(\Gamma_i')$, define $\phi_i'(e) = \phi_i(\psi(e))$ if $\psi(e) \in E(\Gamma_i)$, and $\phi_i'(e) = e$ otherwise. Define $\phi_i''$ similarly. Then each $(\Gamma_i', \phi_i')$ is a $\chi_1$-coloured painting, and each $(\Gamma_i'', \phi_i'')$ is a $\chi_2$-coloured painting, and $\chi_1, \chi_2$ are not bad since $\chi$ satisfies \textbf{S}$_1$. It follows that there exist $j > i \ge 1$ such that there is a linear inflation $\sigma'$ of $(\Gamma_i', \phi_i')$ in $(\Gamma_j', \phi_j')$, and a linear inflation $\sigma''$ of $(\Gamma_i'', \phi_i'')$ in $(\Gamma_j'', \phi_j'')$. \fix{To see that one pair $i, j$ serves both fragments, call a pair $i < j$ \textit{good} if there is a linear inflation of $(\Gamma_i', \phi_i')$ in $(\Gamma_j', \phi_j')$. By Ramsey's theorem there is an infinite set $I \subseteq \{1, 2, \dots\}$ such that either every pair $i < j$ of members of $I$ is good, or none is; and the second is impossible, since then $(\Gamma_i', \phi_i')$ $(i \in I)$ would be a bad sequence for $\chi_1$. Since $\chi_2$ is not bad, the sequence $(\Gamma_i'', \phi_i'')$ $(i \in I)$ is not a bad sequence for $\chi_2$; so there are $i < j$ in $I$ and a linear inflation of $(\Gamma_i'', \phi_i'')$ in $(\Gamma_j'', \phi_j'')$, and the pair $i, j$ is good.}\fixnote{fix:5.5-common-pair}{The paper does not say why one pair $i, j$ serves both fragments. The argument above is the one used in the formalization.} Define $\sigma$ by
\begin{equation*}
\sigma(e) = \begin{cases}
\psi(\sigma'(\psi^{-1}(e))) & \text{if } \psi^{-1}(e) \subseteq \Sigma_1, \\
\psi(\sigma''(\psi^{-1}(e))) & \text{if } \psi^{-1}(e) \subseteq \Sigma_2,
\end{cases}
\end{equation*}
\begin{equation*}
\sigma(v) = \bigcup_{v' \in \psi^{-1}(v) \cap \Sigma_1} \psi(\sigma'(v')) \cup \bigcup_{v'' \in \psi^{-1}(v) \cap \Sigma_2} \psi(\sigma''(v'')).
\end{equation*}
Then $\sigma$ is a linear inflation of $(\Gamma_i, \phi_i)$ in $(\Gamma_j, \phi_j)$, a contradiction, as required. \fix{To see that $\sigma$ is a linear inflation, we use the argument added at the end of the proof of 5.4, with the following changes. Replace $\Sigma_1$ by $\Sigma_1 \cup \Sigma_2$, $\Gamma'_i$ by $\Gamma'_i \cup \Gamma''_i$, $\Gamma'_j$ by $\Gamma'_j \cup \Gamma''_j$, and $\Phi'$ by $\psi^{-1}(\Phi_\chi) = \Phi_{\chi_1} \cup \Phi_{\chi_2}$; let the colour of an edge of $\Gamma'_i \cup \Gamma''_i$ be given by $\phi'_i$ or $\phi''_i$, and similarly for $j$; and replace $\sigma'$ by the map $\tau$ that agrees with $\sigma'$ on the vertices and edges of $\Gamma'_i$, and with $\sigma''$ on those of $\Gamma''_i$. Then $\sigma(e) = \psi(\tau(\psi^{-1}(e)))$ and $V(\sigma(v)) = \bigcup(\psi(V(\tau(v'))) : v' \in \psi^{-1}(v))$, as in 5.4. The facts about $\psi$ listed there still hold. The argument uses only the following properties of $\sigma'$, and $\tau$ has them too. (i) Each edge $f$ is mapped to an edge of the same size, which is internal if and only if $f$ is, borders the same side of $\Phi'$ as $f$ if $f$ is a border edge, has its $k$th end in the image of the $k$th end of $f$ for each $k$, and has a colour at least that of $f$. (ii) Distinct edges have distinct images; the image of each vertex $y$ is a connected induced subdrawing of $sk(\Gamma'_j \cup \Gamma''_j)$, which contains $y$ if $y \in V(\Phi')$; and the images of distinct vertices are disjoint. (iii) Each border edge of $\Gamma'_j \cup \Gamma''_j$ that is not the image of an edge is an edge of the image of some vertex. Indeed, $\sigma'$ and $\sigma''$ have these properties, since they are linear inflations respecting $\Phi_{\chi_1}$ and $\Phi_{\chi_2}$; and so does $\tau$, since $\sigma'$ maps into $\Sigma_1$ and $\sigma''$ maps into $\Sigma_2$, and $\Sigma_1 \cap \Sigma_2 = \emptyset$.}\fixnote{fix:5.5-linear-proof-added}{The paper does not say why $\sigma$ is a linear inflation. In the formalization the two cases, one fragment and two fragments, are treated together in this way.} $\square$

\sectionhead{6. Inflations and linear inflations}

Our strategy to prove 4.1 is to apply a theorem of [10], which we describe later. That, however, applies to inflations rather than linear inflations, and the objective of this section is to smooth over the discrepancy, by means of the following.

\textbf{6.1.} \textit{Let $\chi$ satisfy} \textbf{S}$_1$, \textbf{S}$_3$, \textit{and let} $(\Gamma_i, \phi_i)$ $(i = 1, 2, \dots)$ \textit{be a similarly oriented bad sequence for $\chi$. Then there exists $h \ge 1$ such that for all $j > i \ge h$ there is no inflation of $(\Gamma_i, \phi_i)$ in $(\Gamma_j, \phi_j)$.}

6.1 is a consequence of 5.5 and the following.

\textbf{6.2.} \textit{For any surface $\Sigma$ there is a number $n$ with the following property. Let $\Phi$ be a frame in $\Sigma$, and let $\Gamma_1, \Gamma_2$ be paintings in $\Sigma$ fitting $\Phi$, with $rep(\Gamma_1) \ge n$. Let $\sigma$ be an}
